\clearpage
\section{Box permutation: safety and beta opening}
\label{sec:copying-box-permutation-opening}
This milestone discharges two complete obligations for the concrete BPerm
pre-template: output closure and operational input safety. It also proves the
local application overlap: opening the two related lambda boxes takes one
compute step on each side and produces recursively isomorphic graphs after
compensating the free-interface permutation. Including search, both paths have
exactly six transitions. Their successors have the same exact successful
continuation costs. Bodies may contain nested boxes and arbitrary operations;
the argument retains the original stability premise.

The proof follows the no-term-input criterion in the paper's
\href{https://arxiv.org/html/1907.01257v5\#A8.SS3.SSS2}{Appendix~H.3.2}
and the application-opening case in
\href{https://arxiv.org/html/1907.01257v5\#A8.SS3.SSS6}{Appendix~H.3.6}.
These appendix numbers agree between PDF and HTML. BPerm itself is the literal
Figure~21 pair already constructed earlier. The local diagram is one ingredient
of its Table~3 robustness proof; it does not establish the full specimen-wide
robustness predicate. Copying overlaps, deep overlaps, residual specimen
construction and converse robustness remain to be assembled.

The distinction between external interface permutation and permutation inside
a closed box is essential. The former has a new exact-cost machine theorem,
proved using existing context closure. The latter is still the BPerm template
whose full semantic justification remains unfinished. No new machine rule,
term equation or translation clause is introduced.

\subsection{A typed barrier to entering a template}
\termdef{copying-no-term-inputs}
A pre-template $T$ has \emph{no term inputs} if, whenever $T(G,H)$ holds,
every entry of their common input interface differs from the term wire type.
Such a graph is one way: a term-input-to-term-output operation path would need
a term input, contradicting the interface condition. The left graph alone
therefore witnesses the disjunction required by output closure. The same
reasoning applies to the right graph, although that extra fact is unnecessary.

\lean{Copying/TemplateInputSafety}{10}{11}{PreTemplate.NoTermInputs}
\lean{Copying/TemplateInputSafety}{14}{18}{oneWay_of_no_term_inputs}
\lean{Copying/TemplateInputSafety}{21}{22}{PreTemplate.NoTermInputs.outputClosed}

To use this typing observation operationally, consider a shallow hole of an
actual specimen whose left filling is globally hole free. Its recursive edge
certificate cannot be a retained corolla: the retained edge would still have
the hole label, violating hole freedom. Nor can it be a box certificate, since
a primitive hole label and a box label are disjoint. The certificate must be
a genuine related replacement, hence supplies $T(F_e,K_e)$.

If the specimen's term token entered that hole, its coordinate would occur
on the hole's source list. The ordered source-port typing equation would then
make the corresponding replacement input a term wire. This contradicts the
no-term-input condition. Thus entering is impossible, including when other
replacements are deep inside boxes. Input safety follows by eliminating its
entering premise, for every allowed-context predicate and every pair of cost
relations. There is no assumption of termination or of robustness in this proof.

\lean{Copying/TemplateInputSafety}{26}{37}{DeepFillingEdge.relation_at_hole}
\lean{Copying/TemplateInputSafety}{42}{47}{Specimen.not_entering}
\lean{Copying/TemplateInputSafety}{52}{56}{Specimens.inputSafe_of_no_term_inputs}

\Needspace{10\baselineskip}
The BPerm pre-template relates graphs by the previously defined literal box
permutation up to recursive isomorphism, and requires focus freedom on the left.
Preservation of primitive-label predicates supplies focus freedom on the right.
Destructing a BPerm witness reveals the common input interface
$[\mathsf{thunk}(n)]$. Its sole entry cannot be a term wire. Applying the two
general results proves BPerm output closure and operational input safety.
The latter is quantified over arbitrary $Q,P$, so it also supplies the
$Q^{\mathrm{op}},P$ instance needed by the converse sufficiency theorem.

\lean{Copying/TemplateInputSafety}{60}{62}{BoxPermutation.template}
\lean{Copying/TemplateInputSafety}{64}{75}{BoxPermutation.template_no_term_inputs}
\lean{Copying/TemplateInputSafety}{77}{77}{BoxPermutation.template_outputClosed}
\lean{Copying/TemplateInputSafety}{80}{82}{BoxPermutation.template_inputSafe}

\subsection{Two free-output blocks and the bound-output splice}
For typed permutations $p$ of $F$ and $q$ of $O$, the block action on tagged
indices applies $p$ on the left and $q$ on the right. Conjugating by the
existing split/join bijection gives a typed permutation $p\oplus q$ of $F{+\!+}O$.
The inverse laws follow componentwise. Typing follows from the split/join
label equations and the typing laws of $p,q$. Its two action lemmas state that
joining a left or right index commutes with the corresponding permutation.

\lean{Copying/SpliceBoundary}{9}{14}{TypedPermutation.blocks}
\lean{Copying/SpliceBoundary}{17}{25}{TypedPermutation.append}
\lean{Copying/SpliceBoundary}{27}{29}{TypedPermutation.append_left}
\lean{Copying/SpliceBoundary}{30}{32}{TypedPermutation.append_right}

Let $B$ have input $[\mathsf{term}]$ and outputs $\mathsf{term}::F$, and let
$V$ have input $[\mathsf{term}]$ and outputs $O$. The existing micro-beta splice
connects output zero of $B$ to the input of $V$. Its remaining outputs are
$F{+\!+}O$. The explicit body embedding sends the input to the outer input,
body internals to their tagged summand, bound output zero to the shared wire,
and each free output to the left output block. The argument embedding sends
its input to that same shared wire, its internals to the other tagged summand,
and its outputs to the right output block.

The two raw embeddings are the original simultaneous-expansion insertions.
Case analysis proves them equal to the explicit formulas. The input interface
has exactly one index. For body outputs, the proof splits index zero from a
successor index and uses the original splice pattern's target list. For
argument outputs it uses that pattern's right target block. These equations
ensure that the following comparison is about the literal existing splice.

\lean{Copying/SpliceBoundary}{38}{43}{Application.SpliceBoundary.bodyInsert}
\lean{Copying/SpliceBoundary}{45}{50}{Application.SpliceBoundary.argumentInsert}
\lean{Copying/SpliceBoundary}{52}{54}{Application.SpliceBoundary.rawBodyInsert}
\lean{Copying/SpliceBoundary}{56}{58}{Application.SpliceBoundary.rawArgumentInsert}
\lean{Copying/SpliceBoundary}{60}{72}{Application.SpliceBoundary.body_insert}
\lean{Copying/SpliceBoundary}{74}{87}{Application.SpliceBoundary.argument_insert}

\Needspace{10\baselineskip}
\termdef{copying-splice-action}
A \emph{bound-fixing splice action} consists of a boundary isomorphism
$f:B\to C$ whose output action is $1\oplus p$: the bound output is fixed and
the free tail follows $p$. On spliced vertices, use $f$ on body internals,
identity on the shared wire and argument internals, identity on the outer
input, and $p\oplus 1_O$ on outputs. The displayed inverse uses the inverse
component maps, proving a genuine bijection.

The body-embedding square commutes by three cases. Internal vertices follow
$f$; the sole input is fixed; on outputs the action premise fixes zero and
moves each successor index according to $p$. The argument-embedding square
commutes pointwise, using the right-block action lemma on its outputs.
The edge bijection applies $f$ to body edges and fixes argument edges.

\lean{Copying/SpliceBoundary}{89}{113}{Application.SpliceBoundary.ports}
\lean{Copying/SpliceBoundary}{116}{116}{Application.SpliceBoundary.Action}
\lean{Copying/SpliceBoundary}{118}{141}{Application.SpliceBoundary.body_commutes}
\lean{Copying/SpliceBoundary}{143}{151}{Application.SpliceBoundary.argument_commutes}
\lean{Copying/SpliceBoundary}{153}{159}{Application.SpliceBoundary.edges}

These bijections define a boundary isomorphism between the two splices.
Vertex typing and boundary roles are checked by summands. Labels follow $f$
on body edges and reflexivity on argument edges. For each source or target
list, expand the original insertion formula, apply $f$'s incidence equation,
and use the appropriate commuting square pointwise. This preserves list order.
Finally apply the inverse external action to the target splice. Compensation
fixes every external boundary position, yielding the stronger recursive
isomorphism used by the machine's representation laws.

\lean{Copying/SpliceBoundary}{162}{203}{Application.SpliceBoundary.iso}
\lean{Copying/SpliceBoundary}{206}{211}{Application.SpliceBoundary.compensate}

\subsection{External permutations preserve exact machine costs}
\termdef{copying-permutation-context}
For a symbol $o:I\to O$ and a typed permutation $p$ of $O$, form a
\emph{permutation corolla context} by applying $p$ to the corolla's outputs.
Plugging a graph $G:I\to O$ into its sole edge introduces an empty host-vertex
summand and an impossible surviving-host-edge summand. The existing left-unit
bijections remove those summands. The insertion equation says that this
projection fixes internal vertices and inputs and sends output $i$ to $p(i)$.
Mapping that equation along each incidence list gives a recursive isomorphism
from the filled context to the actual interface-permuted graph.

\lean{Copying/PermutationExecution}{13}{13}{InterfacePermutation.PlugContext.context}
\lean{Copying/PermutationExecution}{15}{25}{InterfacePermutation.PlugContext.insert}
\lean{Copying/PermutationExecution}{28}{44}{InterfacePermutation.PlugContext.iso}

Choose $o$ to be a typed hole with the required interfaces. Any existing
combined-machine rewrite lifts into this corolla by the original
\texttt{Rewrite.plug} constructor. Rename both filled graphs by the proved
corolla isomorphisms. The insertion equation identifies the resulting active
coordinates with their permuted coordinates. This proves rewrite transport
for all operation families at once. Allocation, store access, copying and
beta retain exactly the premises in the original supplied rewrite.

The state action keeps the token's mode and moves its vertex by the interface
action. Search steps transport by the existing boundary-isomorphism search
law; compute steps use the corolla argument just proved. Induction on a finite
run preserves its length. Finality is equivalent because the mode is unchanged
and the interface action preserves the input-vertex predicate.

\lean{Copying/PermutationExecution}{49}{72}{InterfacePermutation.rewrite}
\lean{Copying/PermutationExecution}{75}{76}{InterfacePermutation.state}
\lean{Copying/PermutationExecution}{78}{82}{InterfacePermutation.step}
\lean{Copying/PermutationExecution}{85}{90}{InterfacePermutation.run}
\lean{Copying/PermutationExecution}{92}{96}{InterfacePermutation.final_iff}

\termdef{copying-interface-costs}
Applying $p$ and then $p^{-1}$ gives a state isomorphic to the original state:
the graph comparison is the existing roundtrip isomorphism, modes agree, and
the token equation is the inverse law on outputs and identity elsewhere.
Consequently the \emph{exact successful cost spectrum} is unchanged:
\[
 \mathsf{TerminatesIn}(n,p\cdot s)
 \quad\Longleftrightarrow\quad \mathsf{TerminatesIn}(n,s).
\]
Forward run transport proves one implication. For the other, transport a run
back by $p^{-1}$ and use the roundtrip state isomorphism. Both final states and
the exact number $n$ are preserved. This theorem is about external resource
positions; it does not equate differently permuted closed box bodies.

\lean{Copying/PermutationExecution}{98}{110}{InterfacePermutation.roundtripState}
\lean{Copying/PermutationExecution}{114}{121}{InterfacePermutation.terminatesIn_iff}

\subsection{Moving compensation through the original assembly}
Inverting a block permutation inverts each component; the identity permutation
is its own inverse. Both equations follow from the explicit inverse functions.
In serial assembly, a permutation of the last component's outputs is compatible
with the preceding component because its input action is identity. The serial
boundary isomorphism therefore has identity input action and the prescribed
output action. Factoring it through actual interface permutation gives a
fixed-boundary recursive isomorphism.

\lean{Copying/PermutationAssembly}{9}{10}{TypedPermutation.append_symm}
\lean{Copying/PermutationAssembly}{11}{11}{TypedPermutation.symm_refl}
\lean{Copying/PermutationAssembly}{16}{29}{InterfacePermutation.serial}

For tensor assembly, tensor the two component boundary isomorphisms. Split an
arbitrary boundary index into its tagged component. The existing tensor input
and output action lemmas then identify the input action as identity and the
output action as $p\oplus q$. Factoring this action again gives a recursive
isomorphism. The identity output permutation is itself recursively isomorphic
to the original graph, since it fixes all boundaries. Composing that comparison
on the right component specializes the tensor theorem to a left-only action
while retaining the literal original right component.

\lean{Copying/PermutationAssembly}{32}{60}{InterfacePermutation.tensor}
\lean{Copying/PermutationAssembly}{63}{64}{InterfacePermutation.reflIso}
\lean{Copying/PermutationAssembly}{67}{70}{InterfacePermutation.tensor_left}

\subsection{The actual BPerm and application graphs}
The translated raw lambda body has output interface $E_{k+1,h}$, whereas the
splice exposes $\mathsf{term}::E_{k,h}$. Their equality follows from the
replicate-and-append definition of $E$. Equality casts compose without changing
numerical positions. Define the raw body permutation by casting $1\oplus p$
back across this equality. Its action predicate specifies the output action
of a supplied body boundary isomorphism $f$.

The opening comparison raises $f$ through one representation allowance and
performs the very same interface cast as \texttt{Application.openBody}.
Padding preserves the output action; composing the casts cancels them, so
the opened comparison has precisely the bound-fixing splice action. This
connects the general splice proof to the unchanged application definition.

\lean{Copying/BoxPermutationOpening}{10}{11}{TypedPermutation.cast_trans}
\lean{Copying/BoxPermutationOpening}{17}{18}{BoxPermutation.Opening.interface}
\lean{Copying/BoxPermutationOpening}{21}{26}{BoxPermutation.Opening.bodyPermutation}
\lean{Copying/BoxPermutationOpening}{29}{29}{BoxPermutation.Opening.Action}
\lean{Copying/BoxPermutationOpening}{31}{32}{BoxPermutation.Opening.openIso}
\lean{Copying/BoxPermutationOpening}{34}{42}{BoxPermutation.Opening.open_action}

Casting $f$ instead to the original lambda-box interface supplies the body
action of the canonical BPerm pair with one bound output. The previously proved
box rule then relates the two actual lambda boxes, with inverse compensation
outside the right box. Lift this pair through the lambda corolla, tensor with
the unchanged argument, and then through the application corolla. The serial
and tensor compensation isomorphisms identify the right-hand source as the
literal application of the changed body with the combined external inverse
action. This source comparison is explicitly a structural BPerm derivation,
not an assumed observational equivalence.

\lean{Copying/BoxPermutationOpening}{45}{52}{BoxPermutation.Opening.lambda_pair}
\lean{Copying/BoxPermutationOpening}{57}{68}{BoxPermutation.Opening.source_pair}

Apply the splice theorem to the opened body comparison. The first result is a
boundary isomorphism with output action $p\oplus1_O$. The compensated result
is a boundary-fixing recursive isomorphism. Its edge equation fixes every
argument edge individually. In particular, the proof has not erased an unused
argument or replaced its store-related edges.

\lean{Copying/BoxPermutationOpening}{71}{73}{BoxPermutation.Opening.afterIso}
\lean{Copying/BoxPermutationOpening}{77}{80}{BoxPermutation.Opening.afterCompensated}
\lean{Copying/BoxPermutationOpening}{83}{85}{BoxPermutation.Opening.argument_edge}

\Needspace{16\baselineskip}
\subsection{Equal-cost beta and opening diagrams}
\termdef{copying-beta-permutation-diamond}
Write $L_B(V)$ for the original application redex and $R_B(V)$ for its original
opened splice. Put $q=p\oplus1_O$. For a stable argument $V$, the
\emph{BPerm beta-opening diagram} consists of
\[
\begin{array}{ccc}
 (L_B(V),\mathsf{rewrite},\mathsf{in}_0)&\steps{1}&
 (R_B(V),\mathsf{search},\mathsf{in}_0)\\
 q^{-1}\!\cdot(L_C(V),\mathsf{rewrite},\mathsf{in}_0)&\steps{1}&
 q^{-1}\!\cdot(R_C(V),\mathsf{search},\mathsf{in}_0).
\end{array}
\]
The sources are connected by the structural derivation just proved. The
target graphs are recursively isomorphic, and the resumed input tokens agree.
The upper transition is the existing local beta rule. The lower is that rule
for $C$, transported by the proved external-permutation theorem. The target
comparison is the compensated splice isomorphism. Thus the computation steps
are real combined-machine transitions, not abstract robustness hypotheses.

\lean{Copying/BoxPermutationOpening}{89}{100}{BoxPermutation.Opening.beta_diamond}

The original opening execution has five search transitions followed by beta.
Embedding it into the combined machine gives six transitions on the upper
path; external permutation transports the corresponding six-transition lower
path. The diagram does not assert termination of either body. It is valid
with arbitrary callbacks and recursive box contents, provided the argument
is stable as the original beta rule requires.

\lean{Copying/BoxPermutationOpening}{104}{117}{BoxPermutation.Opening.opening_diamond}

For a literal interface-permuted body $C=p_{\rm body}\cdot B$, the required
body isomorphism and action equation are supplied by the permutation
construction itself. Hence the packaged theorem needs only $B,V,p$ and the
original stability proof: it produces the structural source relation, both
six-step runs and the successor isomorphism. There is no unproved matching,
determinism, robustness or translation-correctness premise hidden in this
instance. The successor state isomorphism further gives equality of all exact
successful continuation costs in either direction. This is a result about
the successors; it is not the full BPerm contextual-equivalence theorem.

\lean{Copying/BoxPermutationOpening}{120}{122}{BoxPermutation.Opening.permuted_action}
\lean{Copying/BoxPermutationOpening}{126}{146}{BoxPermutation.Opening.permuted_opening}
\lean{Copying/BoxPermutationOpening}{149}{157}{BoxPermutation.Opening.successor_costs}

\Needspace{30\baselineskip}
\subsection{Regression witnesses and remaining dependencies}
The regression permutation swaps two free variables and independently swaps
two atoms. Its separate action equations show that neither block acts
trivially. The raw body's bound output still stays at zero, checked directly
from the prefix-extension and cast formulas. The body is the original
translation of a lambda containing an application and an atom occurrence,
with three variables and two atoms in the surrounding scope. Thus opening
the outer application retains a further closed lambda box. Its stable
argument is the literal unit graph at the matching allowance.

\lean{Copying/BoxPermutationOpeningExamples}{10}{11}{BoxPermutation.Opening.Examples.swap}
\lean{Copying/BoxPermutationOpeningExamples}{14}{15}{BoxPermutation.Opening.Examples.nestedBody}
\lean{Copying/BoxPermutationOpeningExamples}{17}{17}{BoxPermutation.Opening.Examples.argument}
\lean{Copying/BoxPermutationOpeningExamples}{19}{21}{BoxPermutation.Opening.Examples.swaps_variable}
\lean{Copying/BoxPermutationOpeningExamples}{23}{25}{BoxPermutation.Opening.Examples.swaps_atom}
\lean{Copying/BoxPermutationOpeningExamples}{28}{31}{BoxPermutation.Opening.Examples.fixes_bound}

Instantiating the packaged diagram proves the compensated nested example's
six-transition execution. Instantiating the successor theorem proves identical
successful continuation costs for every $n$. The edge regression identifies
the surviving unit-argument edge under the comparison. These checks exercise
both resource sorts, bound/free separation, a nested box and the actual
combined machine, without modifying the translation or proposition terms.

\lean{Copying/BoxPermutationOpeningExamples}{34}{40}{BoxPermutation.Opening.Examples.nested_opening}
\lean{Copying/BoxPermutationOpeningExamples}{43}{47}{BoxPermutation.Opening.Examples.nested_successor_costs}
\lean{Copying/BoxPermutationOpeningExamples}{50}{54}{BoxPermutation.Opening.Examples.argument_retained}

Appendix~B.1 removes BPerm output closure and input safety from its pending
obligations, and records this local application diagram as complete. Its
third work package still requires the remaining BPerm overlaps and their
assembly into forward and converse robustness, as well as the other template
families. In particular, duplicating or discarding boxes and transferring
relations beneath nested boxes require further proofs. The focused-graph
correspondence, semantic renaming/substitution, weakening and Figure~22
assembly remain open. Estimated completion remains 60--70\% of mathematical
effort, with low confidence; counting these declarations would overstate how
much of the difficult template work has been discharged.
